\documentclass[11pt,a4paper]{article}

\usepackage[utf8]{inputenc} % allow utf-8 input
\usepackage[T1]{fontenc}    % use 8-bit T1 fonts
\usepackage{hyperref}       % hyperlinks
\usepackage{url}            % simple URL typesetting
\usepackage{booktabs}       % professional-quality tables
\usepackage{amsfonts}       % blackboard math symbols
\usepackage{nicefrac}       % compact symbols for 1/2, etc.
\usepackage{microtype}      % microtypography
\usepackage{xcolor}         % colors
\usepackage{graphicx}       % graphics
\usepackage[normalem]{ulem} % strikethrough
\title{Zeichenblatt-1}

\begin{document}

\maketitle

Blockdiagramm

Sekundärer
Speicher

Eingabegerät

Speichereinheit

Ausgabegerät

Steuereinheit

Arithmetik- und
Logikeinheit

\begin{table}[h]
\begin{tabular}{|l|l|}
\hline
From & To \\ \hline
Eingabegerät & Speichereinheit \\ \hline
Speichereinheit & Ausgabegerät \\ \hline
Steuereinheit & Eingabegerät \\ \hline
Steuereinheit & Ausgabegerät \\ \hline
Steuereinheit &  \\ \hline
 & Arithmetik- und Logikeinheit \\ \hline
 & Arithmetik- und Logikeinheit \\ \hline
\end{tabular}
\end{table}

\end{document}